\section{Complex numbers and the complex plane}

We use $\mathbb{N}$, $\mathbb{Z}$, $\mathbb{R}$, and $\mathbb{C}$ for the natural, integer, real, and complex number systems. A map $f:A\to B$ assigns each $x\in A$ a unique value $f(x)\in B$. It is injective if equal outputs require equal inputs, surjective if every point of $B$ occurs as an output, and bijective if both conditions hold. These ideas identify the complex plane with ordered pairs of real numbers.

A complex number is an ordered pair of real numbers written as
\begin{equation}
z=x+\mathrm{i}y,\qquad x,y\in\mathbb{R},\qquad \mathrm{i}^2=-1.
\end{equation}
The numbers $x=\operatorname{Re}z$ and $y=\operatorname{Im}z$ are respectively its real and imaginary parts. Identifying $z$ with the point $(x,y)$ gives the Argand diagram and turns algebraic operations into plane geometry.

\begin{figure}[htb]
    \centering
    \subimport{tikz/}{complex_plane.tex}
    \caption{Polar representation of a complex number in the Argand plane.}
    \label{fig:ch06-complex-plane}
\end{figure}

In polar form,
\begin{equation}
\rho=|z|=\sqrt{x^2+y^2},\qquad
z=\rho(\cos\varphi+\mathrm{i}\sin\varphi)=\rho e^{\mathrm{i}\varphi}.
\end{equation}
The argument is multivalued: $\operatorname{Arg}z=\arg z+2\pi n$, where $-\pi<\arg z\leq\pi$ is the principal argument. Consequently multiplication and division have the simple forms
\begin{equation}
z_1z_2=\rho_1\rho_2e^{\mathrm{i}(\varphi_1+\varphi_2)},
\qquad
\frac{z_1}{z_2}=\frac{\rho_1}{\rho_2}e^{\mathrm{i}(\varphi_1-\varphi_2)}.
\end{equation}

Addition and multiplication follow from the ordered-pair rules
\begin{equation}
(x_1,y_1)+(x_2,y_2)=(x_1+x_2,y_1+y_2),
\end{equation}
\begin{equation}
(x_1,y_1)(x_2,y_2)=(x_1x_2-y_1y_2,x_1y_2+x_2y_1).
\end{equation}
The additive and multiplicative identities are $0=(0,0)$ and $1=(1,0)$. For $z\ne0$,
\begin{equation}
z^{-1}=\frac{\bar z}{|z|^2}=\frac{x-\mathrm{i}y}{x^2+y^2},
\qquad
\frac{z_1}{z_2}=\frac{x_1x_2+y_1y_2}{x_2^2+y_2^2}
+\mathrm{i}\frac{y_1x_2-x_1y_2}{x_2^2+y_2^2}.
\end{equation}

Complex conjugation is $\bar z=x-\mathrm{i}y$. It reflects a point in the real axis and gives
\begin{equation}
|z|^2=z\bar z,\qquad
\operatorname{Re}z=\frac{z+\bar z}{2},\qquad
\operatorname{Im}z=\frac{z-\bar z}{2\mathrm{i}}.
\end{equation}
The triangle inequality reads $\bigl||z|-|w|\bigr|\leq|z\pm w|\leq|z|+|w|$.
If a polynomial has real coefficients, every nonreal root occurs with its conjugate. De Moivre's formula,
\begin{equation}
(\cos\theta+\mathrm{i}\sin\theta)^n=\cos(n\theta)+\mathrm{i}\sin(n\theta),
\end{equation}
both describes roots and generates multiple-angle identities.

\begin{example}
To solve $z^3=-\mathrm{i}$, write $-\mathrm{i}=e^{\mathrm{i}(3\pi/2+2\pi k)}$. The three distinct roots are
\begin{equation}
z=\mathrm{i},\qquad z=-\frac{\sqrt{3}}{2}-\frac{\mathrm{i}}{2},\qquad
z=\frac{\sqrt{3}}{2}-\frac{\mathrm{i}}{2}.
\end{equation}
They have modulus one and are equally spaced by $2\pi/3$ on the unit circle.
\end{example}

\begin{example}
The locus $\operatorname{Re}(1/z)=2$ obeys $x/(x^2+y^2)=2$, or
\begin{equation}
\left(x-\frac14\right)^2+y^2=\left(\frac14\right)^2.
\end{equation}
It is therefore a circle through the origin with center $(1/4,0)$.
\end{example}

\begin{example}[Polar conversion]
For $z_1=-1+\sqrt3\,\mathrm{i}$, $|z_1|=2$ and $\arg z_1=2\pi/3$, hence
\begin{equation}
z_1=2\left(\cos\frac{2\pi}{3}+\mathrm{i}\sin\frac{2\pi}{3}\right)=2e^{2\pi\mathrm{i}/3}.
\end{equation}
For $z_2=1+\cos\theta+\mathrm{i}\sin\theta$, the half-angle identities give
\begin{equation}
z_2=2\cos\frac\theta2\left(\cos\frac\theta2+\mathrm{i}\sin\frac\theta2\right)
=2\cos\frac\theta2e^{\mathrm{i}\theta/2}.
\end{equation}
\end{example}

Taking the imaginary part of $(\cos\theta+\mathrm{i}\sin\theta)^5$ gives the complete De Moivre calculation
\begin{align}
\sin5\theta
&=5\cos^4\theta\sin\theta-10\cos^2\theta\sin^3\theta+\sin^5\theta\nonumber\\
&=5\sin\theta-20\sin^3\theta+16\sin^5\theta.
\end{align}